% Generated draft: requires mathematical and source review.
\documentclass[11pt]{article}
\usepackage[margin=25mm]{geometry}
\usepackage{fontspec}
\setmainfont{Latin Modern Roman}
\usepackage{amsmath,amssymb,mathtools,unicode-math}
\setmathfont{Latin Modern Math}
\usepackage{enumitem,needspace,xurl,hyperref}
\setlength{\parindent}{0pt}
\setlength{\parskip}{5pt}
\setlength{\emergencystretch}{3em}
\newcommand{\N}{\mathbb{N}}
\newcommand{\Z}{\mathbb{Z}}
\newcommand{\Q}{\mathbb{Q}}
\newcommand{\R}{\mathbb{R}}
\newcommand{\C}{\mathbb{C}}
\newcommand{\F}{\mathbb{F}}
\newcommand{\A}{\mathbb{A}}
\newcommand{\PP}{\mathbb{P}}
\newcommand{\E}{\mathbb{E}}
\newcommand{\eps}{\varepsilon}
\newcommand{\dd}{\,\mathrm{d}}
\begin{document}
\section*{UnsolvedMath ID: 30000228}
\textbf{Source problem code:} catalogue entry
\subsection{Definitions and mathematical statement}
\textbf{Two questions.} Does signature uniqueness extend to rough paths with roughness parameter \(p>1\)? Can one characterize the range of the signature map? The original passage asks both questions; the range question is not removed by a uniqueness theorem.\par In a conventional finite-dimensional continuous weakly geometric formulation, fix \(d\) and \(p>1\), and let \(\mathbf X\) be a based rough path over \(\R^d\), taking values in the step-\(\lfloor p\rfloor\) free nilpotent group and having finite \(p\)-variation with its standard rough-path grading. Its full signature is the sequence\[S(\mathbf X)=\bigl(1,S^1(\mathbf X),S^2(\mathbf X),\ldots\bigr),\qquad S^m(\mathbf X)=\int_{0<t_1<\cdots<t_m<T}d\mathbf X_{t_1}\otimes\cdots\otimes d\mathbf X_{t_m},\] where iterated integrals beyond the defining truncation are interpreted through the rough-path extension. The target is a specified completion of \(\prod_{m\ge0}(\R^d)^{\otimes m}\), not merely the first \(\lfloor p\rfloor\) levels. The source does not fix this category, dimension, completion, or starting-point convention, so this formulation is explicit and conditional.\par The intended uniqueness question is whether equality of full signatures determines the path up to the appropriate tree-like equivalence, rather than pointwise equality of parameterized paths. In the relevant geometric formulation, tree-like equivalence identifies paths when concatenating one with the reversal of the other produces a tree-like excursion, described by a factorization through a real tree with the two endpoints identified; the exact rough-path version must be specified. Increasing reparameterizations and initial translations are already invisible to signatures, so initial points should be fixed or translation quotiented. No assertion that equal truncated signatures determine a rough path is intended.\par For fixed category and parameters, the range question asks for intrinsic conditions on a tensor series \(a=(1,a^1,a^2,\ldots)\) that are necessary and sufficient for \(a=S(\mathbf X)\) for some admissible \(\mathbf X\). The conditions must address both algebraic signature identities and analytic admissibility in the selected tensor topology. Specifying only a formal group-like series, or simply defining the range as this image, is not a supplied characterization. Fixed \(p\) versus a union over roughness parameters, geometric versus weakly geometric paths, and finite-dimensional versus broader ambient spaces are distinct variants left open by the excerpt.\par The supplied triage checked 2026-08-21 labels the record partially solved. It reports the uniqueness subquestion settled through tree-like equivalence and later reconstruction results, while saying a full intrinsic range characterization is not established in the extracted generality. The precise hypotheses of the cited theorems are not reproduced and must be checked before extending that report to a broader path category. This is a dated assessment, not independent theorem verification here.
\subsection{Short English statement}
The two tasks are uniqueness up to tree-like equivalence from the full signature, and intrinsic recognition of tensor series that actually arise as signatures. The supplied triage reports progress settling uniqueness in the cited rough-path setting, while leaving the general range question unresolved.
\subsection{Sources}
\begin{itemize}\item Local UnsolvedMath registry, permanent ID 30000228, code catalogue entry; \emph{Uniqueness and Signature Ranges of Rough Paths}, Oberwolfach Reports Open Problems.\item \emph{Stochastic Analysis and Non-Classical Random Processes (2006)}, Oberwolfach Reports, DOI 10.4171/owr/2005/23, pp. 1231--1298. \url{https://doi.org/10.4171/owr/2005/23}. Proposed year and year tag: 2005; no proposer named.\item Supplied literature label \emph{original Oberwolfach Report 23/2005}: \url{https://doi.org/10.4171/owr/2005/23}.\item Supplied literature label \emph{Boedihardjo--Geng--Lyons--Yang, signature uniqueness}: \url{https://arxiv.org/abs/1406.7871}.\item Supplied literature label \emph{Geng, reconstruction for rough-path signatures}: \url{https://arxiv.org/abs/1508.06890}.\item Literature triage checked 2026-08-21, status \texttt{partially\_solved}; record updated 2026-09-16. The \texttt{corrected\_verified} field reports direct source checking, not independently verified here.\end{itemize}
\end{document}
